\documentclass[11pt, a4paper]{article}
\usepackage[a4paper, margin=2.5cm]{geometry}
\usepackage{fontspec}
\usepackage[spanish, bidi=basic, provide=*]{babel}
\babelprovide[import, onchar=ids fonts]{spanish}
\babelfont{rm}{Noto Sans}
\usepackage{amsmath}
\usepackage{enumitem}

\begin{document}

\section*{Resolución del Ejercicio 171}

\textbf{Enunciado:} Si la probabilidad de que un individuo sufra una reacción por una inyección de un determinado suero, en 0,001, determinar la probabilidad de que de un total de 2.000 individuos:
\begin{enumerate}[label=\alph*)]
    \item exactamente 3,
    \item más de 2 individuos tengan reacción.
\end{enumerate}

\vspace{0.5cm}

\subsection*{1. Planteo del Modelo Original (Binomial)}

Contamos la cantidad de individuos que sufren una reacción alérgica (éxito) dentro de una muestra finita. Como los individuos son independientes y la probabilidad es constante, la variable aleatoria $X$ (número de reacciones) sigue en esencia una distribución Binomial:
\[ X \sim Bi(n = 2000, \, p = 0.001) \]

Si intentáramos calcular esto usando la fórmula binomial estricta, por ejemplo para $k=3$, tendríamos que calcular:
\[ P(X = 3) = \binom{2000}{3} (0.001)^3 (0.999)^{1997} = \frac{2000!}{3! \cdot 1997!} (0.001)^3 (0.999)^{1997} \]
Manejar factoriales de números tan gigantes es muy engorroso (y computacionalmente pesado). 

\subsection*{2. Aproximación a Poisson}

Afortunadamente, existe una regla en probabilidad: cuando el número de ensayos $n$ es muy grande ($n \ge 30$) y la probabilidad de éxito $p$ es muy pequeña ($p \le 0.1$), el cálculo Binomial converge y puede ser aproximado con tremenda precisión usando una \textbf{distribución de Poisson}.
Para ello, el parámetro $\lambda$ (esperanza) de la nueva distribución de Poisson se define igualando el valor esperado de la Binomial:
\[ \lambda = n \cdot p = 2000 \cdot 0.001 = 2 \]

Por ende, trabajaremos con la aproximación:
\[ X \approx \text{Poisson}(\lambda = 2) \]
La fórmula a utilizar será: $P(X = k) = \frac{e^{-\lambda} \lambda^k}{k!} = \frac{e^{-2} 2^k}{k!}$

\vspace{0.5cm}

\subsection*{Parte a) Exactamente 3 individuos tengan reacción}

Nos piden $P(X = 3)$. Reemplazamos directamente en nuestra fórmula de Poisson:
\[ P(X = 3) \approx \frac{e^{-2} \cdot 2^3}{3!} \]
\[ P(X = 3) \approx \frac{e^{-2} \cdot 8}{6} = \frac{4}{3} e^{-2} \]

Sabiendo que $e^{-2} \approx 0.135335$:
\[ P(X = 3) \approx \frac{4}{3} \cdot 0.135335 \approx 0.180447 \]

\textbf{Respuesta a):} La probabilidad de que exactamente 3 individuos tengan reacción es aproximadamente \textbf{0,180}. (Coincide con la guía).

\vspace{0.5cm}

\subsection*{Parte b) Más de 2 individuos tengan reacción}

Aquí se pide la probabilidad de que ocurran estrictamente ``más de 2'' casos, es decir, $P(X > 2)$. 
Lo calculamos por complemento:
\[ P(X > 2) = 1 - P(X \le 2) = 1 - \big[ P(X=0) + P(X=1) + P(X=2) \big] \]

Calculamos las tres probabilidades puntuales necesarias:
\begin{itemize}
    \item $P(X=0) = \frac{e^{-2} \cdot 2^0}{0!} = e^{-2} \approx 0.135335$
    \item $P(X=1) = \frac{e^{-2} \cdot 2^1}{1!} = 2e^{-2} \approx 0.270670$
    \item $P(X=2) = \frac{e^{-2} \cdot 2^2}{2!} = \frac{4e^{-2}}{2} = 2e^{-2} \approx 0.270670$
\end{itemize}

Sumamos estas probabilidades:
\[ P(X \le 2) = e^{-2} + 2e^{-2} + 2e^{-2} = 5e^{-2} \approx 5 \cdot 0.135335 = 0.676675 \]

Finalmente, hallamos el complemento:
\[ P(X > 2) = 1 - 0.676675 = 0.323325 \]

\vspace{0.3cm}
\noindent\fbox{
\parbox{0.95\textwidth}{
\textbf{Nota sobre la respuesta de la guía:} \\
La guía indica que la respuesta al inciso b) es \textbf{0,594}. \\
Ese número, $0.594$, surge de hacer el cálculo para \textbf{``2 o más individuos''} ($P(X \ge 2)$) en lugar de ``más de 2''. 
Si calculamos $P(X \ge 2) = 1 - P(X \le 1) = 1 - (e^{-2} + 2e^{-2}) = 1 - 3e^{-2} = 1 - 3(0.135335) = 1 - 0.4060 = \mathbf{0.594}$. \\
Por lo tanto, la respuesta formal y literal al texto es $0.3233$, pero matemáticamente la guía resolvió el caso $P(X \ge 2)$.
}
}

\end{document}